\section{Local differential geometry of the repaired model (NS-72)}
\label{sec:ns72-local-geometry}

\paragraph{Scope.}
The repaired kernel in NS-70 has a positive local differential-energy
representation. This makes its selected numerator uniformly comparable
to an explicit discrete curvature energy, and supplies two directions
that can be selected in a linear number of scalar operations once the
actual correlation vector is known. These are statements about the
model geometry. Every arithmetic gain below is measured with the
complete actual Gram. No cofinal arithmetic estimate is supplied.

Put
\[
 a_0=K_1-\tfrac32>0,\qquad b_0=K_1/4+\tfrac18=a_0/4+\tfrac12,
 \qquad \alpha_0=\sqrt{b_0/a_0}>\tfrac12.
\]
The normalized stationary kernel is
\[
 \phi_0(r)=e^{-|r|/2}(2K_1+K_1|r|+r^2/4),\qquad
 K_0(u,v)=\frac{\phi_0(\log u-\log v)}{\sqrt{uv}}.
\]
Its complete Fourier density is
\begin{equation}
 w_0(t)=\frac{a_0t^2+b_0}{(t^2+1/4)^3}
       =\frac{a_0}{(t^2+1/4)^2}
                       +\frac{1/2}{(t^2+1/4)^3}.
 \label{eq:ns72-positive-split}
\end{equation}

\begin{proposition}[A causal factor and a local positive energy]
\label{prop:ns72-local-energy}
With Fourier convention $\widehat f(t)=\int f(r)e^{-itr}dr$, define
\[
 h_0(r)=\sqrt{a_0}\,e^{-r/2}
       \left\{r+\frac{\alpha_0-1/2}{2}r^2\right\}
       \mathbf1_{r\geq0}.
\]
Then $h_0\geq0$ and
\[
 \widehat h_0(t)=\frac{\sqrt{a_0}(\alpha_0+it)}{(1/2+it)^3},
 \qquad \phi_0(r)=\int_{\mathbb R}h_0(u+r)h_0(u)\,du.
\]
Let $\mathcal H_0$ be the real Hilbert space with norm
\[
 \|f\|_{\mathcal H_0}^2=
       \frac1{2\pi}\int_{\mathbb R}\frac{|\widehat f(t)|^2}{w_0(t)}\,dt.
\]
As a set it is exactly $H^2(\mathbb R)$, and $\phi_0(r-\cdot)$
is its point-evaluation representer. If $z\in H^3(\mathbb R)$ is the
unique solution of $z'+\alpha_0z=f$, then
\begin{equation}
 \|f\|_{\mathcal H_0}^2
 =\frac1{a_0}\left\{
 \|z'''\|_2^2+\frac34\|z''\|_2^2
               +\frac3{16}\|z'\|_2^2+\frac1{64}\|z\|_2^2\right\}.
 \label{eq:ns72-local-functional}
\end{equation}
In particular, with $L_0=-d^2/dr^2+1/4$,
\begin{equation}
 \frac1{4b_0}\|L_0f\|_2^2\leq\|f\|_{\mathcal H_0}^2
                 \leq\frac1{a_0}\|L_0f\|_2^2.
 \label{eq:ns72-sobolev-comparison}
\end{equation}
All these are whole-line statements; no boundary or remote tail is omitted.
\end{proposition}
\begin{proof}
The elementary transforms of $r e^{-r/2}$ and $r^2e^{-r/2}/2$ on
$r\geq0$ give the displayed causal factor. Its squared modulus is
exactly~\eqref{eq:ns72-positive-split}, proving the autocorrelation
identity. In particular this factorization includes the complete model,
not just its large-frequency term.

The positive rational weight $1/w_0$ is comparable to $(1+t^2)^2$.
Point evaluation is continuous because $\int w_0(t)dt<\infty$;
Fourier inversion gives its stated representer. The multiplier
$(\alpha_0+it)^{-1}$ is an isomorphism from $H^2$ to $H^3$.
It is also the causal convolution with
$e^{-\alpha_0r}\mathbf1_{r\geq0}$; thus there is no unspecified
homogeneous solution. Plancherel gives
$\|f\|_{\mathcal H_0}^2=a_0^{-1}\|(d/dr+1/2)^3z\|_2^2$.
Expanding $(t^2+1/4)^3$ gives~\eqref{eq:ns72-local-functional}.
Equivalently one may use density of compact smooth functions in $H^3$;
no integration-by-parts boundary term survives that limit.
Finally
\[
 \frac1{w_0(t)}=(t^2+1/4)^2\frac{t^2+1/4}{a_0t^2+b_0},
 \qquad \frac1{4b_0}\leq\frac{t^2+1/4}{a_0t^2+b_0}\leq\frac1{a_0},
\]
which proves~\eqref{eq:ns72-sobolev-comparison}.
\end{proof}

\paragraph{The corresponding differential equation.}
For prescribed values $f(r_j)=y_j$, the minimum-norm interpolant is
$f=\sum_j\lambda_j\phi_0(r-r_j)$. In terms of its $z$, it satisfies
in distributions
\begin{equation}
 \frac1{a_0}L_0^3z
     =\sum_j\lambda_j\{\alpha_0\delta_{r_j}-\delta'_{r_j}\}.
 \label{eq:ns72-local-equation}
\end{equation}
Thus between sample points $L_0^3z=0$: its pieces are combinations
of $e^{r/2}$ and $e^{-r/2}$ times quadratic polynomials.
The whole-line Sobolev condition fixes the decaying tails.
The derivatives through order three are continuous, while the jumps are
$[z^{(4)}]_{r_j}=a_0\lambda_j$ and
$[z^{(5)}]_{r_j}=-a_0\alpha_0\lambda_j$.
These signs follow directly from~\eqref{eq:ns72-local-equation}.
This is a positive local model problem with explicit matching conditions;
its coefficients contain no zeta-zero information.

\begin{proposition}[The actual selected numerator is a model interpolation energy]
\label{prop:ns72-interpolation}
Fix $N\geq2$ and the actual optimized correlations $g_n$,
$N<n\leq2N$, from NS-71. Put $r_j=\log j$ and define
\begin{equation}
 y_j=0\quad(j\leq N),\qquad
 y_j=\frac1{\sqrt j}\sum_{k=N+1}^j k g_k\quad(N<j\leq2N).
 \label{eq:ns72-data}
\end{equation}
Then the numerator for the projected repaired model is exactly
\begin{equation}
 m_N=g^TH_{0,N}^{-1}g
       =\min\{\|f\|_{\mathcal H_0}^2:
                          f(r_j)=y_j,\ 1\leq j\leq2N\}.
 \label{eq:ns72-interpolation-numerator}
\end{equation}
This equality holds for every real vector $g$, not just the arithmetic one.
For the actual vector, $y_j=\sqrt j\langle R_N,\sigma_j\rangle$.
\end{proposition}
\begin{proof}
Write $q_j=\langle R_N,\sigma_j\rangle$.
Old-space orthogonality gives $q_j=0$ for $j\leq N$.
The exact difference relation is
$g_j=q_j-(j-1)q_{j-1}/j$; telescoping gives
$q_j=j^{-1}\sum_{k=N+1}^jkg_k$, with no endpoint error.

Let $\Phi_{ij}=\phi_0(r_i-r_j)$ and
$\mathsf D=\operatorname{diag}(\sqrt1,\ldots,\sqrt{2N})$.
The full model Gram is $\mathsf D^{-1}\Phi\mathsf D^{-1}$.
The invertible change from old atoms and new atoms to old atoms and
new differences sends the load data to $(0,g)$.
The complete block-inverse identity therefore gives
$y^T\Phi^{-1}y=g^TH_{0,N}^{-1}g$.
For completeness, the minimum-norm interpolant lies in the span of
the evaluation representers: subtracting its component orthogonal to
that span preserves all values and lowers its squared norm.
Its coefficients are $\Phi^{-1}y$, proving the remaining equality.
Every old/new cross term is retained in this argument.
\end{proof}

For arbitrary data~\eqref{eq:ns72-data}, write
\[
 h_j=r_{j+1}-r_j,\qquad
 d_i=\frac{y_{i+1}-y_i}{h_i}-\frac{y_i-y_{i-1}}{h_{i-1}}
       \quad(2\leq i\leq2N-1).
\]
Let $T$ be the positive tridiagonal matrix indexed by those interior
points, with
\[
 T_{ii}=\frac{h_{i-1}+h_i}{3},\qquad T_{i,i+1}=\frac{h_i}{6}.
\]
Define the complete bending energy and the local curvature sum by
\[
 \mathcal C_N=d^TT^{-1}d,\qquad
 \mathcal W_N=\sum_{i=2}^{2N-1}\frac{d_i^2}{h_{i-1}+h_i}.
\]
They vanish only when $g=0$: zero curvature would make all data affine
in $r$, and the two old zero values force that affine function to vanish.

\begin{proposition}[Uniform comparison with discrete curvature, within the model]
\label{prop:ns72-curvature}
For every $N\geq2$ and every real $g$,
\begin{equation}
 2\mathcal W_N\leq\mathcal C_N\leq6\mathcal W_N,
 \qquad
 \frac{\mathcal C_N}{4b_0}\leq m_N
                         \leq\frac{84}{a_0}\mathcal C_N.
 \label{eq:ns72-curvature-comparison}
\end{equation}
The constants are independent of $N$; the upper constant is deliberately
not optimized. These are uniform comparisons for the model numerator,
not comparisons between the model Gram and the actual arithmetic Gram.
\end{proposition}
\begin{proof}
Let $\psi_i$ be the standard tent function, equal to one at $r_i$
and zero at $r_{i-1},r_{i+1}$, with support between those two endpoints.
Integration by parts on each interval gives
$d_i=\int f''\psi_i$ for any $H^2$ interpolant. The complete Gram of
the tents is exactly $T$. Minimizing $\|f''\|_{L^2(r_1,r_{2N})}^2$
therefore gives $d^TT^{-1}d$. Conversely a minimizing second derivative
in the tent span integrates, with its affine part chosen to match the
first two data, to the natural cubic spline through every sample.
Thus this is an equality for the finite interval, not merely an
uncertified spline approximation. Moreover the off-diagonal row sum
of $T$ is at most half its diagonal. Applying
$2|uv|\leq u^2+v^2$ to its cross terms shows
\[
 \tfrac16\operatorname{diag}(h_{i-1}+h_i)
 \leq T\leq\tfrac12\operatorname{diag}(h_{i-1}+h_i),
\]
and inversion proves the first pair of bounds.
The left side of~\eqref{eq:ns72-sobolev-comparison} bounds every
whole-line interpolant below by $\mathcal C_N/(4b_0)$.

For the uniform upper bound let $F$ be the natural spline on
$[r_1,r_{2N}]$, and set $E=\|F''\|_2^2=\mathcal C_N$ there.
All old nodal values vanish. On an old interval of length $h\leq\log2<1$,
$F'$ has mean zero. Elementary one-dimensional estimates give
$\|F'\|_2^2\leq h^2\|F''\|_2^2$ and
$\|F\|_2^2\leq h^4\|F''\|_2^2$ on that interval.
The derivatives at $r_1$ and $r_N$ have squared magnitudes at most
$E/3$: integrate $F''$ against the endpoint linear weight on an
adjacent zero-data interval. The new interval $[r_N,r_{2N}]$ has
length $\log2<1$, starts at value zero, and hence
\[
 \|F'\|_{L^2(r_N,r_{2N})}^2,\quad
 \|F\|_{L^2(r_N,r_{2N})}^2,\quad
 |F(r_{2N})|^2,\quad |F'(r_{2N})|^2\ \leq\frac83 E.
\]
Combining old and new intervals yields
$\|F'\|_2^2,\|F\|_2^2\leq11E/3$ on the full finite interval.
For any interval $I$ define the positive local functional
\[
 \mathcal J_I(F)=\int_I\left(|F''|^2+\tfrac12|F'|^2
                                      +\tfrac1{16}|F|^2\right)dr.
\]
Consequently $\mathcal J_{[r_1,r_{2N}]}(F)\leq49E/16<4E$.
Only after summing over the complete extension below do we use
$\mathcal J_{\mathbb R}(F)=\|L_0F\|_2^2$. On a finite piece the
two expressions can differ by an endpoint term; none is discarded.

Extend to the left over an interval of length one by the cubic with
zero value and derivative at its outer endpoint and value zero,
derivative $F'(r_1)$ at $r_1$. Its local $\mathcal J$ energy is
$(6833/1680)|F'(r_1)|^2<2E$.
On the right use $y(1-3t^2+2t^3)+w(t-2t^2+t^3)$,
$0\leq t\leq1$, where $y=F(r_{2N})$, $w=F'(r_{2N})$,
and then extend by zero. Its local $\mathcal J$ energy matrix has entries
\[
 A_{00}=7069/560,\quad A_{01}=20339/3360,\quad A_{11}=6833/1680.
\]
They are positive and $A_{00}+2A_{01}+A_{11}<29$.
Since $y^2,w^2,|yw|\leq8E/3$, the right extension costs at most
$232E/3$. The resulting compactly supported function is $C^1$
and belongs to $H^2(\mathbb R)$; jumps in its second derivative
are allowed. Its total $L_0$ energy is less than
$(4+2+232/3)E<84E$.
The upper side of~\eqref{eq:ns72-sobolev-comparison} proves the result.
The extensions explicitly control both tails and all joining conditions.
\end{proof}

\paragraph{Two local coefficient rules.}
Let $\mathsf S_N$ be the exact linear map $g\mapsto d$ defined above,
and put $D_h=\operatorname{diag}(h_{i-1}+h_i)$. The choices
\begin{equation}
 v_{\rm spline}=\mathsf S_N^TT^{-1}\mathsf S_Ng,
 \qquad v_{\rm diagonal}=\mathsf S_N^TD_h^{-1}\mathsf S_Ng
 \label{eq:ns72-local-directions}
\end{equation}
have exact numerators $g^Tv_{\rm spline}=\mathcal C_N$ and
$g^Tv_{\rm diagonal}=\mathcal W_N$.
Both directions are nonzero for $g\ne0$.
Cumulative sums construct $y$, local differences construct $d$, a
positive-pivot tridiagonal solve computes $T^{-1}d$, and weighted suffix
sums apply the transpose. Thus direction selection takes $O(N)$ scalar
operations after the actual $g$ has been supplied. Computing that
arithmetic vector and its true energy is a separate task.

For either direction the actual gain remains
\[
 Q_N(v)=\frac{(g^Tv)^2}{v^TH_Nv}.
\]
A cofinal lower bound for $\mathcal C_N/E_N$ (or $\mathcal W_N/E_N$)
and an upper bound for the corresponding selected true cost, giving a
nonsummable relative contraction, would suffice for RH by the same
argument as NS-71. The geometric equivalence does not establish these
arithmetic bounds. In particular no inequality in
\eqref{eq:ns72-curvature-comparison} includes the actual Gram $H_N$.

\paragraph{Certified finite outcome.}
The following outward intervals enclose each direction's fraction of the
full actual projected gain. Direction selection uses the displayed rule
without fitting its coefficients to the actual Gram.
\[
\begin{array}{c|ccc}
 N&\text{repaired model}&\text{spline curvature}&\text{diagonal curvature}\\ \hline
16&(.8390,.8391)&(.8529,.8530)&(.9308,.9310)\\
32&(.8252,.8254)&(.8334,.8335)&(.9040,.9041)\\
64&(.9577,.9578)&(.9602,.9603)&(.9783,.9785)\\
128&(.9781,.9782)&(.9795,.9796)&(.9852,.9854)\\
256&(.9705,.9707)&(.9725,.9727)&(.9763,.9764)
\end{array}
\]
The simpler diagonal rule outperforms the repaired-kernel solve at all
five checked sizes. This does not assert monotonicity or asymptotic
optimality. The uniform numerator minorant $\mathcal C_N/(4b_0)$ captures
between $6.19\%$ and $6.26\%$ of the exact model numerator on these blocks;
its constant is conservative. Every complete actual Gram, tail and
old-space correction is retained.

Arb runs at 256 and 384 bits match all 140 recorded scalar enclosures.
Doubling the Gram-series cutoff at $N=256$ matches another 28 quantities;
25 overlapping quantities match the published NS-71 calculation.
The minimum recorded accuracy is 73 bits. At $N=16$ an independent
dense solve also encloses all entries of the tridiagonal solution.
The 41 exact Maxima checks include the causal factor, actual Green-function
jumps, normalization, tent Grams and the full extension energies.
No finite computation supplies the missing cofinal arithmetic bound.
